\scanpage{183}
% archive sheet number: 182
Si tout $M\in C$ est complètement réductible,
\inserted{et $C$ défini par un \uline{groupe} $\underline G$ avec une algèbre $A$ sur $G$,}
alors $\Gamma$ s'identifie à l'algèbre des \deleted{\unsure{coéfficients}}
\inserted{\unsure{combinaisons}} \deleted{caractères} linéaires des caractères (fonctions sur
$G$) des représentations \unsure{données|linéaires|degrés|dimensions} de $G$,
$\Gamma^+$ à \deleted{\unsure{un}} \inserted{l'ensemble des} caractères des représentations linéaires
de $G$ (avec conservation de \deleted{\unsure{toutes}} les
structures : $n(f)$ devient $f(e)$ ($e\in G$),
$\vee$ devient $\nu$, $e$ devient la fonction $1$ sur $G$,
\unsure{produits produit|produits produits} ; noter d'ailleurs qu'il
s'agit de fonctions centrales). Si p. ex
$G$ est un groupe algébrique semi-simple complexe,
on obtient une connaissance structurale
concrète de $\Gamma$, qui est l'algèbre des
fonctions régulières sur $T/W$, $T$ un
``tore'' maximal, $W$ le groupe de Weyl
(il faudrait essayer de trouver une
caractérisation intrinsèque simple des
caractères irréductibles, i.e. des $\in I$ ; pour
ceci, il faut mettre sur $T/W$ une
\margin{?} structure de hypergroupe, dont on
doit construire les caractères irréductibles)

Cas de $G=k^*$ ($k$ alg\textsuperscript{t} clos pour
simplifier). Alors $\Gamma(G)\simeq\text{algèbre du groupe }Z$
($I(G)\simeq Z$, avec \unsure{ttes} les structures).
